\documentclass[9pt]{article}
\usepackage[a4paper,margin=1.3cm]{geometry}
\usepackage[spanish]{babel}
\usepackage[utf8]{inputenc}
\usepackage{booktabs}
\usepackage{graphicx}
\usepackage{array}
\usepackage{caption}
\usepackage{hyperref}
\usepackage{enumitem}
\setlist{nosep, leftmargin=1.1em}
\captionsetup{font=scriptsize, labelfont=bf}
\renewcommand{\arraystretch}{1.05}
\pagenumbering{gobble}

\title{\vspace{-1.6em}\textbf{Laboratorio 6 --- Transfer Learning y Fine-Tuning}\vspace{-0.6em}}
\author{CC3092 - Deep Learning y Sistemas Inteligentes}
\date{}

\begin{document}
\maketitle
\vspace{-2.3em}

\section*{1. Investigación: transfer learning en PyTorch}
\footnotesize
\texttt{torchvision.models.vgg16} construye la arquitectura VGG-16; pasarle
\texttt{weights=VGG16\_Weights.IMAGENET1K\_V1} carga los pesos preentrenados en ImageNet, y
\texttt{weights.transforms()} expone el preprocesamiento exacto usado en ese preentrenamiento
(resize 256 + center-crop 224 + normalización con medias/std de ImageNet). El modelo se compone de
\texttt{model.features} (13 capas convolucionales, agrupadas en 5 bloques separados por
\texttt{MaxPool2d}), \texttt{model.avgpool} (\texttt{AdaptiveAvgPool2d} a $7\times7$, lo que permite
usar resoluciones de entrada distintas a $224\times224$) y \texttt{model.classifier} (3 capas
densas $4096\to4096\to1000$, reemplazamos la última por una de 10 salidas). Cada \texttt{nn.Parameter}
tiene un flag \texttt{requires\_grad}; poniéndolo en \texttt{False} sobre todos los parámetros de un
bloque se congela (no recibe gradiente ni se actualiza, y PyTorch evita construir el grafo de
autograd para ese cómputo). \texttt{model.train()}/\texttt{model.eval()} cambian el comportamiento
de \texttt{Dropout}/\texttt{BatchNorm} entre entrenamiento e inferencia; \texttt{torch.no\_grad()}
desactiva la construcción del grafo de autograd (se usa en evaluación y, en este trabajo, también
para forzar el forward de las capas congeladas sin overhead de autograd --- evitarlo resultó
$\sim$4$\times$ más lento en nuestras pruebas sobre el backend \texttt{mps}). Los
\texttt{optimizer.param\_groups} permiten asignar learning rates distintos por capa (backbone con lr
bajo, clasificador con lr alto) al mismo optimizador. \texttt{torchinfo.summary} reporta, vía un
forward simbólico, parámetros y MACs (\emph{multiply-accumulates}) por capa; cada MAC $\approx$ 2
FLOPs, por lo que FLOPs $\approx 2\times$MACs. Los parámetros \emph{totales} son todos los pesos de
la arquitectura; los \emph{entrenables} son el subconjunto con \texttt{requires\_grad=True}
(el optimizador solo actualiza esos). En un servidor con GPU NVIDIA, \texttt{torch.cuda.max\_memory\_%
allocated()} y \texttt{torch.cuda.synchronize()} miden memoria pico y sincronizan el cómputo
asíncrono de la GPU; en esta máquina (Apple Silicon, sin CUDA) se usaron los equivalentes del
backend \texttt{mps}: \texttt{torch.mps.current\_allocated\_memory()} (muestreado por batch para
aproximar el pico) y \texttt{torch.mps.synchronize()}.

\section*{2. Configuración experimental}
\footnotesize
CIFAR-10 (50k train / 10k test, 10 clases balanceadas, imágenes $32\times32\times3$). Split
45{,}000/5{,}000 estratificado para entrenamiento/validación; test oficial reservado para la
evaluación final. Para esta corrida (orientada a validar rápidamente el pipeline completo en una
laptop sin GPU dedicada) se entrenó sobre subconjuntos estratificados pequeños de esos splits
(1{,}500 train / 800 val / 1{,}000 test) y pocos epochs por iteración; el notebook y el código
soportan directamente escalar estos tamaños al dataset completo con más tiempo de cómputo. CNN
desde cero: resolución nativa $32\times32$, normalización con medias/std calculadas sobre
CIFAR-10, augmentation (\texttt{RandomCrop} + \texttt{RandomHorizontalFlip}) solo en entrenamiento.
VGG-16 (feature extractor y fine-tuning): redimensionado a $\mathbf{48\times48}$ (en vez de
$224\times224$, por restricciones de cómputo --- $224\times224$ multiplica los píxeles nativos por
49$\times$, $48\times48$ por 2.25$\times$) y normalización con medias/std de ImageNet. Hardware:
Mac con Apple Silicon; esta corrida en particular se ejecutó sobre CPU (en vez del backend
\texttt{mps}) para evitar contención térmica tras corridas previas.

\section*{3. Resultados de las iteraciones (val.)}
\scriptsize
\begin{center}
\begin{tabular}{p{2.6cm}p{4.6cm}rrrr}
\toprule
\textbf{Arquitectura} & \textbf{Iteración} & \textbf{Params (entr.)} & \textbf{Val Acc} & \textbf{Val F1} & \textbf{t/epoch (s)} \\
\midrule
ITERATIONS_TABLE_ROWS
\bottomrule
\end{tabular}
\end{center}

\section*{4. Comparación de rendimiento y recursos (mejor config. por modelo)}
\scriptsize
\begin{center}
\begin{tabular}{p{2.9cm}rrrrrr}
\toprule
\textbf{Modelo} & \textbf{Params total} & \textbf{Params entr.} & \textbf{FLOPs/img} & \textbf{Latencia (ms)} & \textbf{Epochs$\to$mejor} & \textbf{Test F1} \\
\midrule
COMPARISON_TABLE_ROWS
\bottomrule
\end{tabular}
\end{center}

\vspace{-0.3em}
\noindent\textbf{Experimento de cantidad de datos (10\% vs. 100\% del train):}\\
DATA_QUANTITY_LINE

\begin{center}
\includegraphics[width=0.48\linewidth]{FIG_VAL_ACC_VS_TIME}
\includegraphics[width=0.48\linewidth]{FIG_F1_VS_FLOPS}
\end{center}

\section*{5. Discusión y análisis}
\footnotesize
\textbf{Mejor modelo en test:} BEST_MODEL_DISCUSSION

\textbf{Desempeño vs. recursos:} TRADEOFF_DISCUSSION

\textbf{Features preentrenadas a $32\times32$:} RESIZE_DISCUSSION

\textbf{Fine-tuning (bloques descongelados y lr del backbone):} FINETUNE_DISCUSSION

\textbf{Efecto de la cantidad de datos:} DATA_QUANTITY_DISCUSSION

\textbf{Matrices de confusión:} CONFUSION_DISCUSSION

\textbf{Despliegue (servidor GPU vs. móvil):} DEPLOYMENT_DISCUSSION

\section*{6. Conclusiones}
\footnotesize
CONCLUSIONS_TEXT

\vspace{0.3em}
\noindent\footnotesize\textbf{Repositorio:} \url{REPO_URL}

\end{document}
